\section{Security Gap and Problem Formulation}

\subsection{Authentication Is Not Authorization}

SAGA uses identity checks and access-control tokens to govern contact between
agents. These mechanisms answer different questions. Authentication establishes
which participant completed a protocol step; token validation establishes that
a presented token satisfies the checks performed by the receiving agent.
Authorization additionally asks whether the receiving agent's user-defined
contact policy permits the initiating agent to contact it. A credential can be
authentic and valid under the token checks while the decision to make that
credential available disagrees with the policy's intended meaning.

The reference communication model focuses on participant authentication and
token secrecy. These guarantees concern the origin of protocol messages and
the confidentiality of credentials; they do not, by themselves, establish
that a contact complies with the receiving user's policy. In particular,
confidential delivery to an authenticated participant does not establish that
the participant is entitled to receive a credential under that policy.

The missing requirement is an authorization basis for accepted contact,
preserved through credential issuance and message acceptance. Chapter~4
examines implementation disagreements with policy meaning. This chapter
defines the requirement without asserting an end-to-end access violation.

% Evidence: proofs/proverif/agent_communication.pv, authentication and secrecy queries.

\subsection{Policy--Authorization Consistency}

We call the requirement \emph{policy--authorization consistency}: every
accepted access must be traceable to an applicable authorization decision.
We apply it separately at each authorization boundary: releasing contact
material, issuing a credential, and accepting a message.
For every execution and every access event $e$, we require
\[
  \mathsf{AccessAccepted}(e)
  \Longrightarrow
  \exists d.\;\mathsf{AllowDecision}(d)
  \land \mathsf{Justifies}(d,e).
\]
Here, $d$ is a decision occurring in that execution, and
$\mathsf{AllowDecision}(d)$ means that it grants permission under the intended
policy semantics. An implementation's report of ``allow'' alone does not
establish this predicate. The relation $\mathsf{Justifies}(d,e)$ requires that
the decision precede the accepted access, cover the same subject, target, and
operation, and remain applicable under the relevant policy state. Where a
credential mediates access, the decision must also be linked to the grant and
credential used for that access. An unrelated earlier allow decision is not
a sufficient justification. This is a general requirement whose concrete
events and policy-state conditions must be fixed for each formal model.

Let $I$ be an initiating agent, $R$ a receiving agent, and $P_R$ the contact
rulebook selected by $R$'s user. Each rule associates an agent-identity pattern
with a budget. We write $\mathsf{Allow}_{\mathrm{spec}}(P_R,I)$ for the decision
prescribed by the intended policy semantics, rather than the value returned by
any particular implementation of the matcher. In the cases considered here,
the unique highest-scoring matching rule determines the budget; a positive
budget permits contact, whereas a nonpositive budget or no matching rule does
not. Rule selection follows the matcher's documented most-specific-rule
contract and its AID specificity score; the positive-budget convention follows
the Provider's access check.\footnote{\raggedright See \texttt{match()} and
\texttt{aid\_specificity()} in
\texttt{saga/\allowbreak common/\allowbreak contact\_policy.py}, and
\texttt{access()} in \texttt{saga/\allowbreak provider/\allowbreak provider.py}.}
We initially fix the rulebook throughout an execution and consider either no
matching rule or a unique highest-scoring match. Equal-score conflicts require
an explicit convention and are outside this initial scope.

We distinguish three analysis-level events.
$\mathsf{ProviderGrant}(I,R)$ denotes the Provider's release, through its
access interface, of the material needed for $I$ to continue the contact
protocol with $R$. The event $\mathsf{TokenIssued}(R,I,T)$ denotes $R$'s
creation of token $T$ for $I$ and its release for delivery through the
protocol; it does not assert that $I$ has received or used $T$.
Finally, $\mathsf{MessageAccepted}(R,I,T,m)$ denotes $R$'s acceptance of an
application message $m$ from $I$ with token $T$, after the receiving agent's
access checks, for application processing. It does not assert the completion
of a downstream tool action.

Under this static-policy scope, two necessary conditions at the Provider and
message-acceptance boundaries are
\[
  \mathsf{ProviderGrant}(I,R)
  \Longrightarrow \mathsf{Allow}_{\mathrm{spec}}(P_R,I),
\]
\[
  \mathsf{MessageAccepted}(R,I,T,m)
  \Longrightarrow \mathsf{Allow}_{\mathrm{spec}}(P_R,I).
\]
The analogous requirement at token issuance is
\[
  \mathsf{TokenIssued}(R,I,T)
  \Longrightarrow \mathsf{Allow}_{\mathrm{spec}}(P_R,I).
\]
All three implications range over every occurrence of the respective events
and all their arguments. They express permission at specific enforcement
boundaries; they do not replace the decision-to-access linkage required by
$\mathsf{Justifies}$. Nor does policy permission alone suffice for access:
certificate checks, key availability, token validity, expiry, and quota may
still prevent a permitted contact. These conditions address contact
permission, not quantitative bounds on credential issuance or message use.

For a policy that can change over time, the applicable policy version and the
effect of revocation on existing credentials must additionally be specified.
Quota consumption and token expiry likewise require explicit state and
validity conditions. The general requirement leaves room for these
obligations, but the static permission conditions above do not establish them.

A message-level authorization claim therefore requires a model connecting
policy evaluation, credential issuance, and the final message-acceptance event.
Chapter~6 develops the corresponding formalization and states its assumptions.

% Evidence: saga/common/contact_policy.py, match() documentation and default budget.
% Evidence: saga/provider/provider.py, access() decision and OTK response.
% Evidence: saga/agent.py, token checks and receive_conversation().

\subsection{Authorization Pipeline}

The relevant decision is distributed across several stages. The user's
rulebook determines an intended answer for $I$ contacting $R$. The Provider
evaluates its stored copy of that rulebook and, if the request passes its
checks, returns one-time key material. The initiating and receiving agents
then complete the credential exchange; the receiving agent generates a token
and later validates a token presented with a message. An authorization
argument must connect the same initiating identity and receiving identity
across these stages, rather than infer permission solely from the presence of
a valid token.
The intended progression is
\[
  P_R \longrightarrow \mathsf{ProviderGrant}
  \longrightarrow \mathsf{TokenIssued}
  \longrightarrow \mathsf{MessageAccepted}.
\]
The arrows indicate protocol stages and authorization dependencies, not a
guarantee that reaching one stage causes the next to succeed.

Each boundary requires separate evidence: a Provider response does not
establish token issuance, and token issuance does not establish message
acceptance. Later checks and protocol steps may fail.

When $\neg\mathsf{Allow}_{\mathrm{spec}}(P_R,I)$ holds, we distinguish three
levels of violation:
\begin{enumerate}
  \item \textbf{Provider decision violation:}
    $\mathsf{ProviderGrant}(I,R)$ occurs despite the policy denial.
  \item \textbf{Credential issuance violation:}
    $\mathsf{TokenIssued}(R,I,T)$ occurs despite the policy denial.
    Issuance alone does not establish delivery or subsequent use.
  \item \textbf{End-to-end authorization violation:}
    $\mathsf{MessageAccepted}(R,I,T,m)$ occurs despite the policy denial.
    This is acceptance at the receiving agent's application boundary.
\end{enumerate}
These levels classify possible observations rather than assert results.
A sound pipeline must preserve the link to the applicable decision,
participants, and credential across stages. Chapter~6 specifies the
corresponding events; Chapter~7 evaluates implementation observations and
model results separately.

% Evidence: saga/provider/provider.py, access() and returned key material.
% Evidence: saga/agent.py, policy check, token issuance, and token validation.

\subsection{Threat Model}

We consider a normally registered initiating agent controlled by the adversary.
It can choose its contact requests and participate in the protocol under its
genuine identity and certificate. The receiving user's rulebook is an existing
configuration: the adversary cannot change its contents or rule order, control
the Provider or CA, or obtain another honest party's private keys.

The Provider and receiving agent execute their implementations with the same
fixed rulebook. Cryptographic checks are assumed effective. The failure under
study is a disagreement between policy meaning and enforcement during ordinary
protocol operation, without compromise of either enforcement point.

The adversary's objective is an execution in which the intended policy denies
contact by $I$, yet the system performs one of the authorization-relevant
actions above. At the message level, the violation condition is
\[
  \mathsf{MessageAccepted}(R,I,T,m)
  \land \neg\mathsf{Allow}_{\mathrm{spec}}(P_R,I).
\]
This end-to-end violation is distinct from an erroneous Provider grant or
credential issuance. Claiming it requires evidence that the intervening
checks and protocol steps reach the receiving agent's acceptance boundary.

We study this objective under a static policy. Policy updates and revocation,
disagreement between stored policy versions, quantitative exhaustion of
budgets or quotas, compromise of trusted parties, and arbitrary actions after
contact require additional assumptions and are outside this initial question.
The distinction between zero, negative, and positive configured budgets
remains part of the contact-permission semantics.
